% Yeban User Guide — Chapter 10: Terminal and Automation Entry Points (English)
% T-DOC-001 chapter-by-chapter delivery, written ahead of the CH10 task.
% Sources: README.md, docs/README.md, docs/tool-availability.md,
% presence of crates/yeban-tui, cli.rs switches. Tool counts are never hand-copied.

\chapter{Terminal and Automation Entry Points}
\label{ch:entry-points}

\classquote{天下同归而殊涂，一致而百虑。}{Book of Changes, “Appended Judgments II” (《周易·系辞下》)}{Translation: “All beneath heaven returns to one end though the roads differ; all reaches one goal though thoughts are many.” The GUI, the terminal, and MCP are different entries into the same project.}





The same project can be operated through three entry points: the desktop GUI, a
terminal TUI and the command line CLI, plus an MCP tool surface for AI agents. They
share one project model and backend, so they do the same work with different
interaction styles. This chapter gives one section each and says when to use which.

\section{TUI}

Yeban provides a terminal user interface (TUI) entry point, with source in
\texttt{crates/yeban-tui}. It serves setups without a graphical desktop, or users who
prefer working in a terminal, without opening a window; the panels and keys it actually
offers are documented by that entry point's own help, not hand-copied here where they
could go stale.

\begin{table}[htbp]
\centering
\caption{What each entry point suits}
\label{tab:ch10-entry-compare}
\begin{tabular}{ll}
\hline
Entry & Best for \\
\hline
Desktop GUI & Everyday arranging, notes, mixing, visual feedback \\
TUI & Terminal environments, no graphical desktop, remote shells \\
CLI & Batch work, scripts, export, CI \\
MCP & AI agents editing and querying through tool calls \\
\hline
\end{tabular}
\end{table}

\section{CLI}

\subsection{Read the help first}

The complete switch list and exit codes come from \texttt{yeban-app --help}; this
chapter is a guide, not a replacement. Common entries:

\begin{description}
  \item[\texttt{--open <path>}] Open a given project directly (GUI or batch).
  \item[\texttt{--export-wav}] Export the full mix; pair with \texttt{--loudness-target} and \texttt{--true-peak-ceiling} for mastering. An unreachable target fails loudly rather than silently degrading.
  \item[\texttt{--export-stems <dir>}] Export stems with the same backend as the MCP tool; track-selection switches can limit it to named tracks.
  \item[\texttt{--apply-ops}] Commit a file of edit operations as one batch.
  \item[\texttt{--export-song-overview}] Export the whole-song overview (energy/frequency/density/emotion with an anomaly summary).
  \item[\texttt{--print-shortcuts}] Print the shortcut policy table (the same source as the GUI's Help $\to$ Shortcuts).
  \item[\texttt{--headless}] Run on a machine without a display.
\end{description}

\subsection{Know the feature gates}

Some switches are gated by build features. In a default build the MCP-related
\texttt{in-process-mcp} is on — but "compiled in" is not "listening": at runtime you
still need \texttt{--enable-mcp-http} or the environment variable
\texttt{YEBAN\_MCP\_HTTP=1}. With \texttt{--no-default-features}, calling a gated
switch yields a usage error naming the feature (exit code 2), never a silent ignore.
Experimental exports such as \texttt{--export-als} and \texttt{--export-logic} are off
by default and only compile when requested.

\section{MCP}

\subsection{Two forms, one registry}

MCP has two forms sharing a single \texttt{TOOLS} registry and the same domain
implementation:

\begin{description}
  \item[Form A (library form)] Hosted in-process by \texttt{yeban-app}, for an AI agent working on a running DAW; transport is loopback HTTP (\textbf{only} \texttt{127.0.0.1}) or direct in-process calls.
  \item[Form B (binary form)] A standalone \texttt{yeban-mcp} binary over stdio for batch work without a GUI.
\end{description}

Network listening binds only the loopback address, never \texttt{0.0.0.0} — a safety
red line. \texttt{yeban-mcp}'s own \texttt{mcp-http} is never in its default features;
only the host can forward-enable it.

\screenshot{ch10-mcp-modes}{The two MCP forms: in-app mount over loopback HTTP, standalone binary over stdio}

\begin{figure}[htbp]
\centering
\begin{tikzpicture}[node distance=1.2cm, >=stealth,
    box/.style={draw, rounded corners, align=center, minimum width=2.4cm, minimum height=0.9cm}]
  \node[box] (a) {One TOOLS registry};
  \node[box, below left=of a] (b) {Form A: in-app mount};
  \node[box, below right=of a] (c) {Form B: standalone binary};
  \node[box, below=of b] (d) {Loopback HTTP 127.0.0.1};
  \node[box, below=of c] (e) {stdio batch};
  \draw[->] (a)--(b); \draw[->] (a)--(c); \draw[->] (b)--(d); \draw[->] (c)--(e);
\end{tikzpicture}
\caption{Both MCP forms share one registry; network listening binds only loopback}
\label{fig:ch10-mcp-forms}
\end{figure}

\subsection{Don't trust hand-copied tool counts}

The tool count must be read live from the three constants in
\texttt{crates/yeban-mcp/src/tools.rs} and the enum in
\texttt{schemas/mcp-tools.schema.json}: numbers in documentation have gone stale
repeatedly and old manuals disagree with each other. This chapter deliberately gives
no concrete number.

Some tools depend on a provider injected by the host (e.g.\ an external SVS singing
service). Form B has no injector, so calling such a tool returns the
\texttt{providerUnavailable} contract rather than pretending success.

\section{Help Topic Mapping}

\begin{table}[htbp]
\centering
\caption{Sections of this chapter $\to$ in-app help anchors}
\label{tab:ch10-help}
\begin{tabular}{ll}
\hline
Section & Help topic \\
\hline
TUI & \texttt{entry-tui} \\
CLI: help / common switches & \texttt{entry-cli-switches} \\
CLI: feature gates & \texttt{entry-cli-features} \\
MCP: two forms & \texttt{entry-mcp-forms} \\
MCP: tool counts and providers & \texttt{entry-mcp-tools} \\
\hline
\end{tabular}
\end{table}

\paragraph{What you will be able to do} Pick the right entry for the situation, know common CLI switches and feature gates, and explain that both MCP forms share one registry while network listening binds only loopback.

\paragraph{Hands-on exercise} Use \texttt{yeban-app --help} to find the exact stem-export switch, export stems from both the GUI and CLI for the same project and compare receipts; run \texttt{yeban-app --print-shortcuts} against the GUI panel; finally confirm that without \texttt{--enable-mcp-http} the MCP HTTP port is not listening.
